% =====================================================================
%  acknowledgments.tex — מקורות (עמוד 19) + על הספר: קרדיטים ואיך נבנה הספר (עמוד 20)
% =====================================================================
\chapter*{מקורות}
\addcontentsline{toc}{chapter}{מקורות}
\markboth{מקורות}{}
\thispagestyle{plain}

הטענות העובדתיות במדריך מבוססות על המקורות שלהלן. הבדיקה של כל טענה מול המקור שלה מתועדת
במסמך ביקורת המקורות (\en{Source Audit}) שבמאגר הפרויקט.

\begin{english}
\printbibliography[heading=none]
\end{english}
\clearpage

% ---------------------------------------------------------------------
\chapter*{על הספר}
\addcontentsline{toc}{chapter}{על הספר}
\markboth{על הספר}{}
\thispagestyle{plain}

\secnotoc{קרדיטים לחומרים חזותיים}
כל האיורים, הגרפים והתרשימים בספר מקוריים ונוצרו בקוד. קבצי המקור שלהם נמצאים בתיקייה \en{figures}
שבמאגר הפרויקט, וכולם נוצרים מחדש בפקודת בנייה אחת.
\begin{itemize}
  \item \textbf{איור השער ואיור \ref{fig:elevator}}: איורים שנבנו בקוד \en{TikZ} מצורות גאומטריות.
  \item \textbf{איור \ref{fig:ladder}}: סולם החשיפה, תרשים בקוד \en{TikZ}.
  \item \textbf{איור \ref{fig:map}}: מפה קונספטואלית שנוצרה בקוד \en{Python}.
  \item \textbf{איור \ref{fig:graph}}: גרף שנוצר בקוד \en{Python} מנתוני המחשה.
  \item \textbf{איור \ref{fig:flow}}: תרשים זרימה בן שמונה שלבים בקוד \en{TikZ}.
\end{itemize}

\secnotoc{איך נבנה הספר}
הספר נכתב ונבנה כפרויקט \en{Vibe Coding}, כלומר פיתוח בהנחיה בשפה טבעית בעזרת \en{Claude Code}.
הטקסט נכתב בעזרת \en{XeLaTeX} עם \en{polyglossia}, לתמיכה בעברית ובכיווניות מימין לשמאל, והגרפים
והאיורים נוצרו בקוד \en{Python} ובשפת השרטוט \en{TikZ}.

תהליך העבודה תוכנן סביב שלושה סוכנים בעלי תפקידים נפרדים, שעובדים על אותו מאגר קבצים:

\begin{center}
\begin{tikzpicture}[font=\sffamily\small,
    agent/.style={rtlpar,draw,rounded corners=3mm,line width=1pt,align=center,text width=3.3cm,
                  minimum height=1.5cm,inner sep=4pt},
    arr/.style={-{Stealth[length=6pt]},line width=0.9pt,draw=navy},
    lbl/.style={font=\sffamily\scriptsize,text=charcoal,fill=paper,inner sep=1.5pt}]
  \node[agent,fill=navy!8!white,draw=navy]          (M) at (11,0)  {\textbf{\en{Master}}\\ כתיבה, עיצוב וניהול};
  \node[agent,fill=coral!14!white,draw=coraldark]   (D) at (5.5,0) {\textbf{\en{Detector}}\\ בדיקת איכות עצמאית};
  \node[agent,fill=sage!20!white,draw=sagedark]     (F) at (0,0)   {\textbf{\en{Fixer}}\\ אימות מקורות ותיקונים};
  \draw[arr] ([yshift=5pt]M.west) -- node[lbl,above=2pt]{\RLE{בקשת בדיקה}} ([yshift=5pt]D.east);
  \draw[arr] ([yshift=-5pt]D.east) -- node[lbl,below=2pt]{\RLE{דוח ליקויים}} ([yshift=-5pt]M.west);
  \draw[arr] (M.north) -- ++(0,0.55) -| node[lbl,pos=0.25,above=1pt]{\RLE{בקשת תיקון ואימות מקורות}} (F.north);
  \draw[arr] (F.south) -- ++(0,-0.55) -| node[lbl,pos=0.25,below=1pt]{\RLE{בקשה לבדיקה חוזרת}} (D.south);
\end{tikzpicture}
\end{center}

\en{Master} כותב ומעצב את הספר ומחליט על קבלת השינויים, \en{Detector} בודק באופן עצמאי את השפה,
הפריסה והעמידה בדרישות, ו\en{Fixer} מאמת את המקורות ומתקן ליקויים שאושרו. מצב הביצוע בפועל של
כל שלב, כולל אילו בדיקות כבר הושלמו, מתועד ביומני התהליך ובדוחות הבדיקה שבמאגר הפרויקט.
בינה מלאכותית שימשה ככלי עבודה, והאחריות לתוכן נשארת של המחברות.

\vfill
\begin{knowbox}[מגישות]
ליה פורת ורוני ברק \enspace·\enspace קוד קבוצה: \en{liyaroni-ex2} \enspace·\enspace 2026
\end{knowbox}
